\documentclass[aspectratio=169,11pt]{beamer}
\usepackage{fontspec}
\usepackage{babel}
\usepackage{booktabs}
\usepackage{xcolor}
\usepackage{graphicx}

\definecolor{coral}{HTML}{cc785c}
\definecolor{ink}{HTML}{141413}
\definecolor{cream}{HTML}{faf9f5}
\definecolor{muted}{HTML}{6c6a64}

\babelprovide[import, onchar=ids fonts]{thai}
\babelprovide[import, onchar=ids fonts]{english}

\babelfont{rm}{Noto Sans}
\babelfont[thai]{rm}{Noto Sans Thai}
\babelfont{tt}{Noto Sans Mono}
\babelfont[thai]{tt}{Noto Sans Thai}

\directlua{
  local glyph_id = node.id("glyph")
  local penalty_id = node.id("penalty")

  local function is_thai(c)
    return c >= 0x0E01 and c <= 0x0E5B
  end

  local function is_thai_combining(c)
    return (c >= 0x0E31 and c <= 0x0E3A) or (c >= 0x0E47 and c <= 0x0E4E)
  end

  luatexbase.add_to_callback("pre_linebreak_filter", function(head)
    for n in node.traverse_id(glyph_id, head) do
      local prev = n.prev
      if prev and prev.id == glyph_id and is_thai(prev.char)
         and is_thai(n.char) and not is_thai_combining(n.char) then
        local p = node.new(penalty_id)
        p.penalty = 0
        node.insert_before(head, n, p)
      end
    end
    return head
  end, "thai_break")
}

\emergencystretch=1em

%% ─── Beamer theme ─────────────────────────────────────────────
\usetheme{default}
\usecolortheme{default}

\setbeamercolor{frametitle}{fg=ink, bg=cream}
\setbeamercolor{title}{fg=ink}
\setbeamercolor{normal text}{fg=ink}
\setbeamercolor{structure}{fg=coral}
\setbeamercolor{block title}{bg=coral!20, fg=ink}
\setbeamercolor{block body}{bg=coral!5}
\setbeamercolor{footline}{fg=muted}
\setbeamercolor{background canvas}{bg=cream}

\setbeamertemplate{navigation symbols}{}
\setbeamertemplate{frametitle}[default][left]
\setbeamerfont{frametitle}{size=\large, series=\bfseries}

\setbeamertemplate{footline}{%
  \hfill\color{muted}\small\insertframenumber/\inserttotalframenumber\hspace{8pt}\vspace{4pt}
}

\setbeamertemplate{itemize item}{\color{coral}\textbullet}
\setbeamertemplate{itemize subitem}{\color{muted}--}

%% ─── Title info ───────────────────────────────────────────────
\title{\textbf{ชื่องานนำเสนอ}}
\subtitle{Subtitle in English}
\author{ชื่อผู้นำเสนอ}
\institute{มหาวิทยาลัย · ภาควิชา}
\date{\today}

\begin{document}

%% ─── Title slide ──────────────────────────────────────────────
\begin{frame}
  \titlepage
\end{frame}

%% ─── Outline ──────────────────────────────────────────────────
\begin{frame}{สารบัญ}
  \tableofcontents
\end{frame}

%% ─── Section 1 ────────────────────────────────────────────────
\section{บทนำ}

\begin{frame}{ปัญหาและแรงจูงใจ}
  \begin{columns}
    \begin{column}{0.6\textwidth}
      \begin{itemize}
        \item ประเด็นที่ 1
        \item ประเด็นที่ 2
        \item ประเด็นที่ 3
      \end{itemize}
    \end{column}
    \begin{column}{0.4\textwidth}
      \begin{block}{จุดสำคัญ}
        ข้อความเน้น
      \end{block}
    \end{column}
  \end{columns}
\end{frame}

%% ─── Section 2 ────────────────────────────────────────────────
\section{วิธีการ}

\begin{frame}{แนวทางที่นำเสนอ}
  \begin{enumerate}
    \item ขั้นตอนที่ 1
    \item ขั้นตอนที่ 2
    \item ขั้นตอนที่ 3
  \end{enumerate}

  \vspace{12pt}
  \begin{alertblock}{สิ่งที่แตกต่างจากงานก่อนหน้า}
    อธิบายความแตกต่างหลัก
  \end{alertblock}
\end{frame}

%% ─── Section 3 ────────────────────────────────────────────────
\section{ผลการทดลอง}

\begin{frame}{ผลลัพธ์หลัก}
  \centering
  \begin{tabular}{@{}lccc@{}}
    \toprule
    \textbf{Method} & \textbf{Metric 1} & \textbf{Metric 2} & \textbf{Metric 3} \\
    \midrule
    Baseline A  & 00.0 & 00.0 & 00.0 \\
    Baseline B  & 00.0 & 00.0 & 00.0 \\
    \textbf{Ours} & \textbf{00.0} & \textbf{00.0} & \textbf{00.0} \\
    \bottomrule
  \end{tabular}
\end{frame}

%% ─── Section 4 ────────────────────────────────────────────────
\section{สรุป}

\begin{frame}{สรุปและงานต่อไป}
  \begin{columns}
    \begin{column}{0.5\textwidth}
      \textbf{สรุปผล}
      \begin{itemize}
        \item ผลสรุปที่ 1
        \item ผลสรุปที่ 2
      \end{itemize}
    \end{column}
    \begin{column}{0.5\textwidth}
      \textbf{งานต่อไป (Future Work)}
      \begin{itemize}
        \item ทิศทางที่ 1
        \item ทิศทางที่ 2
      \end{itemize}
    \end{column}
  \end{columns}
\end{frame}

%% ─── Thank you ────────────────────────────────────────────────
\begin{frame}
  \centering
  \vspace{1cm}
  {\Huge \textbf{ขอบคุณครับ/ค่ะ}}\\[12pt]
  {\large Thank You}\\[24pt]
  {\color{muted} ชื่อ · อีเมล}
\end{frame}

\end{document}
